\documentclass{article}
\usepackage{amsmath, amssymb, amsthm}
\usepackage{mathrsfs}

\title{The Infinite CAR Algebra}
\author{DeepSeek Chat}

\begin{document}

\maketitle

\section{Definition of the Infinite CAR Algebra}
The \textbf{infinite CAR algebra} is the canonical anticommutation relations (CAR) algebra defined over an infinite-dimensional separable Hilbert space. It is a \( C^* \)-algebra generated by the fermionic creation and annihilation operators satisfying the CAR.

\subsection*{Definition:}
Let \(\mathcal{H}\) be an infinite-dimensional separable Hilbert space. The \textbf{infinite CAR algebra}, denoted \(\text{CAR}(\mathcal{H})\), is the universal \( C^* \)-algebra generated by the set \(\{a(f) \;|\; f \in \mathcal{H}\}\) (called \textbf{annihilation operators}) subject to the relations:
\begin{enumerate}
    \item \textbf{Antilinearity in the argument}:  
    \[ a(\lambda f + \mu g) = \overline{\lambda} a(f) + \overline{\mu} a(g) \quad \forall f, g \in \mathcal{H}, \lambda, \mu \in \mathbb{C}. \]
    \item \textbf{Canonical Anticommutation Relations (CAR)}:  
    \[
    \{a(f), a(g)\} = 0 \quad \text{and} \quad \{a(f), a^*(g)\} = \langle f, g \rangle \cdot \mathbf{1},
    \]
    where:
    \begin{itemize}
        \item \(a^*(f)\) is the adjoint of \(a(f)\) (the \textbf{creation operator}),
        \item \(\{A, B\} = AB + BA\) is the anticommutator,
        \item \(\langle \cdot, \cdot \rangle\) is the inner product on \(\mathcal{H}\),
        \item \(\mathbf{1}\) is the unit of the algebra.
    \end{itemize}
\end{enumerate}

\section{Key Properties}
\begin{enumerate}
    \item \textbf{Unique \( C^* \)-Completion}:  
    The CAR algebra is the unique \( C^* \)-algebra (up to \(*\)-isomorphism) satisfying these relations.
    
    \item \textbf{Fermionic Fock Space Representation}:  
    The standard representation is on the \textbf{antisymmetric Fock space} \(\mathcal{F}_a(\mathcal{H})\), where:
    \begin{itemize}
        \item \(a(f)\) annihilates a fermion in state \(f\),
        \item \(a^*(f)\) creates a fermion in state \(f\).
    \end{itemize}
    
    \item \textbf{Irreducibility}:  
    The representation is irreducible if \(\mathcal{H}\) is infinite-dimensional.
    
    \item \textbf{Approximation by Finite CAR Algebras}:  
    If \(\mathcal{H}_n \subset \mathcal{H}\) is an increasing sequence of finite-dimensional subspaces with \(\bigcup_n \mathcal{H}_n\) dense in \(\mathcal{H}\), then:
    \[
    \text{CAR}(\mathcal{H}) \cong \overline{\bigcup_n \text{CAR}(\mathcal{H}_n)},
    \]
    where each \(\text{CAR}(\mathcal{H}_n)\) is isomorphic to a full matrix algebra \(M_{2^n}(\mathbb{C})\).
    
    \item \textbf{\(\mathbb{Z}_2\)-Grading (Even/Odd Operators)}:  
    The CAR algebra has a natural \(\mathbb{Z}_2\)-grading, distinguishing \textbf{even} and \textbf{odd} operators based on the number of creation/annihilation operators.
    
    \item \textbf{Connection to Clifford Algebras}:  
    For a real Hilbert space \(\mathcal{H}_\mathbb{R}\), the CAR algebra is isomorphic to the \textbf{Clifford algebra} \(\text{Cliff}(\mathcal{H}_\mathbb{R} \oplus i \mathcal{H}_\mathbb{R})\).
\end{enumerate}

\section{Is the Infinite CAR Algebra a \( W^* \)-Algebra?}
The \textbf{infinite CAR algebra} is primarily considered as a \textbf{\( C^* \)-algebra}, not a \textbf{\( W^* \)-algebra} (also known as a von Neumann algebra). However, it has representations that generate von Neumann algebras. Below is a detailed explanation:

\subsection*{1. CAR Algebra as a \( C^* \)-Algebra}
\begin{itemize}
    \item The \textbf{infinite CAR algebra}, \(\text{CAR}(\mathcal{H})\), is the universal \( C^* \)-algebra generated by creation and annihilation operators satisfying the canonical anticommutation relations (CAR).
    \item It is constructed as the norm closure of the algebra of finite products of creation and annihilation operators.
    \item As a \( C^* \)-algebra, it is \textbf{not} a von Neumann algebra because:
    \begin{itemize}
        \item It is defined abstractly via generators and relations, not as an algebra of bounded operators on a Hilbert space.
        \item It lacks a natural predual (a necessary condition for being a \( W^* \)-algebra).
    \end{itemize}
\end{itemize}

\subsection*{2. The CAR Algebra in a Fock Representation Yields a von Neumann Algebra}
When represented on a Hilbert space (e.g., the \textbf{fermionic Fock space} \(\mathcal{F}_a(\mathcal{H})\)), the weak closure of \(\text{CAR}(\mathcal{H})\) in \(\mathcal{B}(\mathcal{F}_a(\mathcal{H}))\) \textbf{is} a von Neumann algebra. Specifically:
\begin{itemize}
    \item The \textbf{Fock representation} \(\pi: \text{CAR}(\mathcal{H}) \to \mathcal{B}(\mathcal{F}_a(\mathcal{H}))\) is irreducible if \(\mathcal{H}\) is infinite-dimensional.
    \item The von Neumann algebra generated by this representation is \(\mathcal{B}(\mathcal{F}_a(\mathcal{H}))\) (the algebra of all bounded operators on Fock space), which is a \textbf{type I factor}.
    \item If \(\mathcal{H}\) is finite-dimensional, the weak closure is a finite-dimensional matrix algebra (a type I factor).
\end{itemize}

\subsection*{3. Other Representations and von Neumann Algebras}
Different representations of the CAR algebra can lead to different von Neumann algebras:
\begin{itemize}
    \item \textbf{Quasi-free representations} (e.g., the Araki-Wyss representations in quantum statistical mechanics) generate von Neumann algebras that can be \textbf{type III factors} (common in quantum field theory).
    \item The \textbf{GNS representation} associated with a tracial state (for finite \(\mathcal{H}\)) gives a hyperfinite type II\(_1\) factor.
\end{itemize}

\subsection*{4. Key Distinction}
\begin{itemize}
    \item The \textbf{abstract CAR algebra} is a \( C^* \)-algebra.
    \item Its \textbf{concrete representations} generate von Neumann algebras (depending on the representation).
\end{itemize}

\subsection*{Conclusion}
The \textbf{infinite CAR algebra itself} is not a \( W^* \)-algebra, but its representations on Hilbert spaces generate von Neumann algebras (often factors of type I, II, or III). This distinction is crucial in mathematical physics, where the CAR algebra is used abstractly in algebraic quantum field theory, while its von Neumann algebraic versions appear in statistical mechanics and quantum field theory.

\end{document}